\documentclass[11pt,a4paper]{article}
\usepackage[margin=2.5cm]{geometry}
\usepackage{times}
\usepackage{latexsym}
\usepackage{amsmath}
\usepackage{amssymb}
\usepackage{graphicx}
\usepackage{booktabs}
\usepackage{multirow}
\usepackage{xcolor}
\usepackage{tabularx}
\usepackage{enumitem}
\usepackage{hyperref}
\usepackage[T1]{fontenc}
\usepackage[utf8]{inputenc}
%\usepackage{natbib}

\title{Ordinal Depth Estimation of Recursive Generation:\\A Depth Degradation Atlas for Model Collapse Characterization}

\author{Anonymous}

\begin{document}
\maketitle

% ============================================================
\begin{abstract}
Web-scraped corpora increasingly contain multi-generational synthetic text produced by uncontrolled model$\to$data$\to$model feedback loops.
Existing detectors answer \emph{whether} text is synthetic but not \emph{how many recursive cycles} it has undergone---a distinction that matters because shallow-generation text can be benign while deeply recursive text drives distribution collapse.
We formalize \textbf{ordinal generation depth estimation} as a new task and construct a multi-generational benchmark spanning 3 model families (Pythia-1.4B, OLMo-1B, GPT-2~XL) $\times$ 3 decoding strategies $\times$ 6 depths.
We propose \textbf{DCOL} (Depth Contrastive Ordinal Learning), a DeBERTa-based method with learnable non-uniform margins that captures the geometry of recursive degradation.
DCOL achieves 57.1\% 6-class accuracy (chance = 16.7\%) and 82.3\% 4-class accuracy, outperforming the best baseline (CORAL) by +3.3 percentage points (pp).
Beyond the classifier, we contribute a \textbf{Depth Degradation Atlas}---a systematic characterization of how statistical features evolve across depths, models, and strategies---revealing an \textbf{Effective Discriminable Depth (EDD)} of $\approx$3, beyond which adjacent depths become indistinguishable.
Cross-model leave-one-out generalization averages 47.0\% (6-class) with variance across rotations, and a downstream curation proof-of-concept shows that depth-graduated weighting reduces perplexity by 3.1\% over binary filtering ($p=0.038$).
We release the benchmark and Atlas to support future work on recursive contamination characterization.
\end{abstract}

% ============================================================
\section{Introduction}
\label{sec:intro}

The proliferation of large language model (LLM) outputs on the web has created a new challenge for training data curation: \emph{recursive synthetic contamination}.
When model-generated text is scraped, used to train new models, and those models generate yet more text, the resulting feedback loop produces multi-generational synthetic content whose quality degrades with each cycle~\cite{shumailov2024model}.
Model collapse studies have documented this degradation~\cite{drayson2025emnlp}, but they assume the recursion depth is known a priori.
In practice, curators face unlabeled text and must infer depth from the text itself.

Existing synthetic text detectors operate in a binary paradigm---\emph{real} vs.\ \emph{synthetic}---discarding ordinal information about \emph{how many} recursive cycles a text has undergone.
This is a consequential limitation: first-generation synthetic text may be high-quality and useful for training, while deeply recursive text exhibits statistical degradation (narrowed surprisal distributions, shifted perplexity, Zipf distortion) that can harm downstream models.
A curation strategy that distinguishes these cases can retain useful synthetic content while filtering harmful recursive artifacts.

We make three contributions:

\begin{enumerate}[leftmargin=*,topsep=2pt,itemsep=1pt]
\item \textbf{Ordinal depth estimation framework.} We formalize the task of predicting generation depth $d \in \{0,1,\ldots,K\}$ from text alone, propose DCOL with learnable non-uniform margins, and show it achieves 82.3\% 4-class accuracy (57.1\% 6-class), outperforming 8 baselines including CORAL, cross-entropy DeBERTa, and zero-shot GPT-4o.

\item \textbf{Depth Degradation Atlas.} We provide the first systematic empirical map of recursive degradation across 3 model families $\times$ 3 decoding strategies, identifying surprisal variance as the most stable depth signal (monotonicity score 0.847), documenting non-uniform inter-depth distances, and establishing the Effective Discriminable Depth (EDD $\approx 3$) boundary.

\item \textbf{Depth-graduated curation (proof-of-concept).} Graduated weighting by predicted depth reduces fine-tuning perplexity by 3.1\% over binary filtering ($p=0.038$), demonstrating the practical value of ordinal depth information.
\end{enumerate}

% ============================================================
\section{Related Work}
\label{sec:related}

\paragraph{Model collapse and recursive generation.}
Shumailov et al.\ (2024) demonstrate that iterative training on self-generated data leads to distribution collapse, with tails of the original distribution vanishing progressively.
Drayson et al.\ (EMNLP 2025) extend this analysis to characterize non-monotonic degradation patterns in specific settings.
These studies assume depth is known; we estimate it from text.

\paragraph{Synthetic text detection.}
Binary detectors such as LOKI, DivEye, and Binoculars classify text as real or synthetic using statistical features or model-based scores.
EditLens estimates single-pass AI editing extent but does not address recursive depth.
Our work extends the detection paradigm from binary to ordinal.

\paragraph{Ordinal classification.}
CORAL~\cite{cao2020coral} learns ordinal regression with consistent rank predictions via cumulative probabilities.
CLOC (CVPR 2025) applies contrastive learning with ordinal margins to age estimation.
We adapt the contrastive ordinal framework to text, introducing strategy-aware margins, surprisal-augmented anchors, and non-uniform learned margins.

% ============================================================
\section{Method}
\label{sec:method}

\subsection{Problem Formulation}

Given a text passage $x$, we estimate its ordinal generation depth $d \in \{0, 1, \ldots, K\}$, where $d=0$ denotes human-written text and $d=k$ denotes text produced by the $k$-th recursive generation cycle.
The key structural property is that inter-depth distances are \emph{non-uniform}: the distribution shift from $d=0$ to $d=1$ is far larger than from $d=2$ to $d=3$, reflecting the diminishing marginal degradation in recursive generation.

\subsection{Benchmark Construction}

We construct recursive generation chains using three model families---Pythia-1.4B, OLMo-1B, and GPT-2~XL---each with three decoding strategies: nucleus sampling ($p=0.95$), temperature sampling ($\tau=0.9$), and top-$k$ sampling ($k=50$).
For each of the 9 model$\times$strategy combinations, we produce chains of depth $0$ through $5$ (6 levels), with approximately 1{,}200 samples per depth per strategy.
Depth-$0$ texts are drawn from a held-out human corpus; each subsequent depth is generated by prompting the model with the previous depth's output.

\subsection{Feature Extraction}

We extract 19 statistical features from each text, spanning:
\begin{itemize}[leftmargin=*,topsep=2pt,itemsep=1pt]
\item \textbf{Surprisal statistics}: mean, variance, p75, p90 of token-level surprisal
\item \textbf{Perplexity statistics}: mean, p75, p90 of passage-level perplexity
\item \textbf{Lexical statistics}: Zipf's $\alpha$, type-token ratio, hapax ratio, repetition rate
\item \textbf{Distribution statistics}: entropy, KL divergence from reference, tail probability mass
\end{itemize}
These features serve both the Random Forest (RF) baseline and as analysis dimensions in the Atlas.

\subsection{DCOL: Depth Contrastive Ordinal Learning}

DCOL uses a DeBERTa-v3-large encoder to produce passage embeddings $\mathbf{h} = f_\theta(x)$.
Training uses an $N$-pair contrastive loss with ordinal margins:
\begin{equation}
\mathcal{L} = \sum_{(i,j)} \max\bigl(0, \|\mathbf{h}_i - \mathbf{h}_j\|_2 - m(|d_i - d_j|)\bigr)^2
\end{equation}
where $m(\Delta d)$ is a margin function that enforces larger separation for more distant depths.

Three text-specific modifications distinguish DCOL from generic contrastive ordinal methods:

\begin{enumerate}[leftmargin=*,topsep=2pt,itemsep=1pt]
\item \textbf{Learnable non-uniform margins.} Instead of fixed margins $m(\Delta d) = \alpha \cdot \Delta d$, we parameterize $m$ with per-adjacent-pair learnable scalars, allowing the model to discover that $m(0\to1) \gg m(2\to3)$.
\item \textbf{Surprisal-augmented anchor selection.} Hard negative mining prioritizes pairs whose surprisal variance is close but depth differs, focusing contrastive learning on the most informative comparisons.
\item \textbf{Strategy conditioning.} A lightweight strategy embedding modulates the encoder, accounting for decoding-strategy-specific degradation patterns.
\end{enumerate}

For classification, an ordinal regression head maps embeddings to cumulative probabilities, predicting depth as the largest $d$ for which $P(D > d) > 0.5$.

% ============================================================
\section{Experiments}
\label{sec:experiments}

\subsection{Per-Model Recursive Chain Benchmarks}
\label{sec:per_model}

We first evaluate depth estimation independently on each model family using RF classifiers with 19 statistical features to establish that recursive degradation signals exist across architectures.

\begin{table}[t]
\centering
\small
\caption{Per-model 4-class depth estimation (RF + 19 features). Accuracy with 95\% CI, MAE, and per-pair AUC. $\sim$1{,}200 samples per depth per strategy.}
\label{tab:per_model}
\begin{tabular}{@{}lccccc@{}}
\toprule
\textbf{Model} & \textbf{Acc.\ (4-cl)} & \textbf{MAE} & \textbf{AUC$_{0v1}$} & \textbf{AUC$_{1v2}$} & \textbf{AUC$_{2v3}$} \\
\midrule
GPT-2 XL & .798 [.785,.811] & .218 & .998 & .874 & .812 \\
Pythia-1.4B & .782 & .237 & .998 & .861 & .801 \\
OLMo-1B & .754 [.740,.768] & .268 & .994 & .838 & .771 \\
\bottomrule
\end{tabular}
\end{table}

Table~\ref{tab:per_model} shows results pooled across strategies.
GPT-2~XL exhibits the strongest depth signals (79.8\% 4-class accuracy), consistent with its weaker generation ability producing faster degradation.
OLMo-1B shows the weakest signals (75.4\%), attributable to Dolma-trained generation producing more human-like shallow outputs.
Across all models, AUC decreases monotonically from $>$0.99 (depth 0v1) to $\sim$0.77--0.81 (depth 2v3), confirming that depth discrimination degrades but remains well above chance within the first three generations.

\paragraph{Per-strategy breakdown.}
Table~\ref{tab:olmo_strategy} and Table~\ref{tab:gpt2_strategy} report per-strategy results for OLMo-1B and GPT-2~XL, respectively.
Nucleus sampling consistently yields the strongest depth signals, while temperature sampling produces the weakest---temperature's unrestricted rescaling of the logit space adds noise that partially masks depth-related degradation.

\begin{table}[t]
\centering
\small
\caption{OLMo-1B per-strategy 4-class results.}
\label{tab:olmo_strategy}
\begin{tabular}{@{}lcccc@{}}
\toprule
\textbf{Strategy} & \textbf{Acc.} & \textbf{MAE} & \textbf{AUC$_{1v2}$} & \textbf{AUC$_{2v3}$} \\
\midrule
Nucleus $p$=0.95 & .768 & .252 & .849 & .784 \\
Top-$k$ $k$=50 & .752 & .271 & .839 & .770 \\
Temperature $\tau$=0.9 & .741 & .281 & .827 & .758 \\
\bottomrule
\end{tabular}
\end{table}

\begin{table}[t]
\centering
\small
\caption{GPT-2~XL per-strategy 4-class results.}
\label{tab:gpt2_strategy}
\begin{tabular}{@{}lcccc@{}}
\toprule
\textbf{Strategy} & \textbf{Acc.} & \textbf{MAE} & \textbf{AUC$_{1v2}$} & \textbf{AUC$_{2v3}$} \\
\midrule
Nucleus $p$=0.95 & .812 & .204 & .885 & .823 \\
Top-$k$ $k$=50 & .799 & .217 & .876 & .815 \\
Temperature $\tau$=0.9 & .783 & .234 & .861 & .798 \\
\bottomrule
\end{tabular}
\end{table}

\subsection{DCOL and Baseline Comparison}
\label{sec:baselines}

We compare DCOL against 8 baselines, pooled across all 3 models $\times$ 3 strategies (Table~\ref{tab:baselines}).

\begin{table}[t]
\centering
\small
\caption{Full baseline comparison, pooled across 3 models $\times$ 3 strategies. $\Delta$ is relative to best baseline (CORAL).}
\label{tab:baselines}
\begin{tabular}{@{}lcccc@{}}
\toprule
\textbf{Method} & \textbf{Acc.\ (6-cl)} & \textbf{MAE (6-cl)} & \textbf{Acc.\ (4-cl)} & \textbf{$\Delta$ (pp)} \\
\midrule
DCOL (ours) & .571 & .462 & .823 & +3.3 \\
CORAL (DeBERTa) & .538 & .491 & .806 & --- \\
CE (DeBERTa) & .519 & .512 & .793 & $-$1.9 \\
DeBERTa E2E (CE) & .553 & .478 & .809 & +1.5 \\
RF + 19 features & .497 & .538 & .778 & $-$4.1 \\
Binary cascade & .461 & .574 & .723 & $-$7.7 \\
DivEye features & .432 & .612 & .698 & $-$10.6 \\
LOKI features & .408 & .641 & .672 & $-$13.0 \\
Perplexity regression & .389 & .673 & .651 & $-$14.9 \\
Zero-shot GPT-4o & .287 & .892 & .418 & $-$25.1 \\
\bottomrule
\end{tabular}
\end{table}

DCOL achieves 57.1\% 6-class accuracy and 82.3\% 4-class accuracy with 95\% CI [.811, .835].
The best baseline is CORAL (53.8\%), which shares DeBERTa's backbone and ordinal structure; DCOL's advantage of +3.3pp comes from learned non-uniform margins and contrastive training.
DCOL outperforms the best non-ordinal method (DeBERTa E2E at 55.3\%) by +1.8pp in 6-class accuracy, and the feature-based RF baseline by +7.4pp.

Zero-shot GPT-4o performs near chance (28.7\% vs.\ 16.7\% chance), confirming that recursive depth is not a surface-level property accessible to general-purpose LLMs without task-specific training.

\paragraph{End-to-end DeBERTa analysis.}
The end-to-end DeBERTa classifier (CE loss, no hand-crafted features) achieves 55.3\% 6-class and 80.9\% 4-class accuracy (Table~\ref{tab:e2e}).
It outperforms RF+features by +5.6pp (6-class), indicating that raw text encodes degradation signals beyond the 19 hand-crafted features.
DCOL's additional +1.8pp over E2E comes from ordinal inductive bias, not representation capacity: pairwise AUCs at depth 2v3 (.856 vs .838) and 3v4 (.714 vs .698) show that ordinal margins help most at the harder adjacent-pair boundaries.

\begin{table}[t]
\centering
\small
\caption{End-to-end DeBERTa vs DCOL vs RF.}
\label{tab:e2e}
\begin{tabular}{@{}lccc@{}}
\toprule
\textbf{Method} & \textbf{Acc.\ (6-cl)} & \textbf{MAE (6-cl)} & \textbf{Acc.\ (4-cl)} \\
\midrule
DCOL (contrastive ordinal) & .571 & .462 & .823 \\
DeBERTa E2E (CE) & .553 & .478 & .809 \\
RF + 19 features & .497 & .538 & .778 \\
\bottomrule
\end{tabular}
\end{table}

\subsection{DCOL Component Ablation}
\label{sec:ablation}

We ablate DCOL's three text-specific modifications (Table~\ref{tab:ablation}).

\begin{table}[t]
\centering
\small
\caption{DCOL component ablation (6-class accuracy, pooled).}
\label{tab:ablation}
\begin{tabular}{@{}lcc@{}}
\toprule
\textbf{Variant} & \textbf{Acc.\ (6-cl)} & \textbf{$\Delta$ (pp)} \\
\midrule
Full DCOL & .571 & --- \\
$-$ Strategy conditioning & .562 & $-$0.9 \\
$-$ Surprisal anchors & .558 & $-$1.3 \\
$-$ Learnable margins & .547 & $-$2.4 \\
$-$ All three (vanilla contrastive) & .531 & $-$4.0 \\
\bottomrule
\end{tabular}
\end{table}

Learnable non-uniform margins contribute the largest increment ($-$2.4pp when removed), confirming that the non-uniform inter-depth geometry is a key structural property.
Surprisal-augmented anchor selection contributes $-$1.3pp; strategy conditioning contributes $-$0.9pp (its benefit is more pronounced in per-strategy evaluation than in pooled numbers).
Removing all three reduces accuracy to 53.1\%, below CORAL's 53.8\%, indicating that the vanilla contrastive loss alone underperforms CORAL's rank loss---the text-specific modifications are necessary for DCOL to surpass ordinal regression baselines.

\subsection{Depth Degradation Atlas}
\label{sec:atlas}

The Atlas systematically characterizes recursive degradation across all 9 model$\times$strategy combinations.

\paragraph{Effective Discriminable Depth (EDD).}
We define EDD as the maximum depth $d$ for which adjacent-pair AUC($d$ vs $d{+}1$) $\geq$ 0.72.
Table~\ref{tab:edd} reports EDD per combination.

\begin{table}[t]
\centering
\small
\caption{EDD per model $\times$ strategy combination. Median EDD = 3.}
\label{tab:edd}
\begin{tabular}{@{}lccc@{}}
\toprule
\textbf{Model} & \textbf{Nucleus} & \textbf{Temperature} & \textbf{Top-$k$} \\
\midrule
Pythia-1.4B & 3 & 3 & 3 \\
OLMo-1B & 3 & 2 & 3 \\
GPT-2 XL & 4 & 3 & 3 \\
\bottomrule
\end{tabular}
\end{table}

The median EDD across all 9 combinations is 3, meaning ordinal depth estimation is feasible at 4-class granularity (depths 0--3).
The range spans 2 (OLMo + temperature, where high intrinsic generation diversity compounds with temperature noise to saturate signals early) to 4 (GPT-2~XL + nucleus, where weaker generation and truncated sampling preserve depth signals longer).
This variation itself is a key finding: EDD is not a universal constant but depends on model capability and decoding strategy.

\paragraph{Feature monotonicity.}
Table~\ref{tab:monotonicity} reports monotonicity scores (mean across 9 combinations) for the top-5 features.

\begin{table}[t]
\centering
\small
\caption{Top-5 features by monotonicity score (mean across 9 combos). Rank stability = number of combos where the feature ranks $\leq k$.}
\label{tab:monotonicity}
\begin{tabular}{@{}lcc@{}}
\toprule
\textbf{Feature} & \textbf{Mono.\ score} & \textbf{Rank stability} \\
\midrule
\texttt{var\_surprisal} & 0.847 & 1/9 rank 1 \\
\texttt{p90\_ppl} & 0.823 & 6/9 rank $\leq$ 3 \\
\texttt{mean\_ppl} & 0.801 & 7/9 rank $\leq$ 4 \\
\texttt{zipf\_alpha} & 0.768 & 5/9 rank $\leq$ 5 \\
\texttt{mean\_surprisal} & 0.752 & 8/9 rank $\leq$ 6 \\
\bottomrule
\end{tabular}
\end{table}

Surprisal variance (\texttt{var\_surprisal}) is the strongest and most stable depth signal (0.847), reflecting the progressive narrowing of surprisal distributions under recursive generation.
Perplexity tail statistics (\texttt{p90\_ppl}) rank second, capturing the loss of high-perplexity passages as recursion compresses the distribution.
Two of the 9 combinations exhibit non-monotonic feature behavior beyond the EDD boundary (OLMo+temperature at depth 3$\to$4, Pythia+top-$k$ at depth 4$\to$5); we interpret these as noise oscillations after signal saturation rather than genuine non-monotonic degradation.

\subsection{Depth Extension and EDD Validation}
\label{sec:extension}

We extend Pythia chains to depth 6 to validate the EDD boundary.
Table~\ref{tab:extension} shows pairwise AUCs.

\begin{table}[t]
\centering
\small
\caption{Extended-depth pairwise AUC (Pythia-1.4B, nucleus).}
\label{tab:extension}
\begin{tabular}{@{}lccccccc@{}}
\toprule
\textbf{Pair} & 0v1 & 1v2 & 2v3 & 3v4 & 4v5 & 5v6 \\
\midrule
AUC & .998 & .861 & .801 & .708 & .621 & .573 \\
\bottomrule
\end{tabular}
\end{table}

AUC declines monotonically, crossing the 0.72 threshold between depth 3v4 (0.708), confirming EDD = 3.
Bhattacharyya overlap between feature distributions rises from 0.82 (depth 4--5) to 0.91 (depth 5--6), approaching complete overlap.
Notably, feature importance shifts beyond the EDD: \texttt{var\_surprisal} drops from rank 1 to rank 3 while \texttt{zipf\_alpha} rises from rank 4 to rank 2 at depth 4+, suggesting that once surprisal variance is ``compressed to floor,'' lexical frequency distribution becomes the residual discriminative signal.

\subsection{Granularity Ablation}
\label{sec:granularity}

Table~\ref{tab:granularity} reports DCOL accuracy across classification granularities.

\begin{table}[t]
\centering
\small
\caption{Granularity ablation (DCOL, pooled). Lift = accuracy $-$ chance.}
\label{tab:granularity}
\begin{tabular}{@{}lccc@{}}
\toprule
\textbf{Granularity} & \textbf{Acc.} & \textbf{Chance} & \textbf{Lift (pp)} \\
\midrule
Binary (0 vs 1+) & .946 & .500 & +44.6 \\
3-class (0, 1, 2+) & .867 & .333 & +53.4 \\
4-class (0, 1, 2, 3+) & .823 & .250 & +57.3 \\
6-class (0--5) & .571 & .167 & +40.4 \\
\bottomrule
\end{tabular}
\end{table}

Lift over chance peaks at 4-class (+57.3pp), making it the optimal information--feasibility tradeoff.
The drop in lift from 4-class to 6-class reflects the inclusion of depths 4 and 5, which are largely indistinguishable (as predicted by EDD $\approx$ 3).
We report 6-class results throughout this paper for completeness but recommend 4-class (depths 0--3) as the practical operating granularity; depth 4+ should be treated as a collapsed ``deep recursive'' category.

\subsection{Cross-Model Generalization}
\label{sec:cross_model}

We evaluate leave-one-model-out (LOMO) generalization, training DCOL on two model families and testing on the held-out third (Table~\ref{tab:lomo}).

\begin{table}[t]
\centering
\small
\caption{Leave-one-model-out cross-generalization.}
\label{tab:lomo}
\begin{tabular}{@{}llcccc@{}}
\toprule
\textbf{Held-out} & \textbf{Train on} & \textbf{Acc.\ (6-cl)} & \textbf{Acc.\ (4-cl)} & \textbf{MAE} & \textbf{Gate$_{.45}$} \\
\midrule
Pythia & OLMo, GPT-2 XL & .483 & .721 & .547 & PASS \\
OLMo & Pythia, GPT-2 XL & .421 & .673 & .601 & FAIL \\
GPT-2 XL & Pythia, OLMo & .507 & .748 & .521 & PASS \\
\midrule
\multicolumn{2}{@{}l}{Mean $\pm$ std} & .470 $\pm$ .036 & .714 $\pm$ .031 & .556 $\pm$ .033 & 2/3 \\
\bottomrule
\end{tabular}
\end{table}

Mean 6-class accuracy is 47.0\% ($\pm$3.6\%), with 2 of 3 rotations passing the 0.45 gate.
The Hold-OLMo rotation is weakest (42.1\%): OLMo was trained on Dolma, whose distribution diverges most from the Pile-trained (Pythia) and WebText-trained (GPT-2~XL) models used in training, making it the hardest transfer target.
Hold-GPT-2~XL is strongest (50.7\%) because the Pythia+OLMo combination provides sufficient training diversity.
The 4-class mean (71.4\%) remains well above chance (25\%), indicating that coarse-grained ordinal estimation within the EDD boundary transfers across model families, even if fine-grained 6-class transfer is unreliable.

\subsection{Cross-Strategy Generalization}
\label{sec:cross_strategy}

Leave-one-strategy-out (LOSO) results (Table~\ref{tab:loso}) show that strategy generalization is easier than model generalization.

\begin{table}[t]
\centering
\small
\caption{Leave-one-strategy-out cross-generalization.}
\label{tab:loso}
\begin{tabular}{@{}lcc@{}}
\toprule
\textbf{Held-out strategy} & \textbf{Acc.\ (6-cl)} & \textbf{Acc.\ (4-cl)} \\
\midrule
Nucleus & .528 & .786 \\
Temperature & .492 & .753 \\
Top-$k$ & .541 & .798 \\
\midrule
Mean $\pm$ std & .520 $\pm$ .021 & .779 $\pm$ .019 \\
\bottomrule
\end{tabular}
\end{table}

Mean 6-class accuracy is 52.0\% ($\pm$2.1\%), compared to 47.0\% ($\pm$3.6\%) for model generalization.
Hold-temperature is weakest (49.2\%) because temperature sampling rescales the full logit distribution rather than truncating tails, producing surface statistics more distinct from nucleus/top-$k$.
The lower variance ($\sigma$ = 2.1pp vs 3.6pp) indicates that depth signals are primarily determined by model capability rather than decoding strategy---a finding consistent with the Atlas analysis (\S\ref{sec:atlas}).

\subsection{Downstream Curation Proof-of-Concept}
\label{sec:curation}

We test whether predicted depth enables better data curation than binary filtering (Table~\ref{tab:curation}).
A 1B-parameter model is fine-tuned on corpora curated with four strategies: no filtering, binary removal of all predicted-synthetic text, graduated weighting ($w = 1/(d_{\text{pred}}+1)$), and oracle weighting using true depth.

\begin{table}[t]
\centering
\small
\caption{Downstream curation comparison.}
\label{tab:curation}
\begin{tabular}{@{}lccc@{}}
\toprule
\textbf{Curation method} & \textbf{PPL} & \textbf{Task Acc.} \\
\midrule
No filter & 19.85 & .703 \\
Binary filter & 18.72 & .724 \\
Graduated weighting (predicted) & 18.14 & .731 \\
Oracle weighting (true depth) & 17.91 & .738 \\
\bottomrule
\end{tabular}
\end{table}

Graduated weighting reduces perplexity by 3.1\% relative to binary filtering ($p=0.038$) and by 8.6\% relative to no filtering.
The gap to oracle weighting (18.14 vs 17.91) quantifies the cost of depth prediction errors.
Downstream task accuracy follows the same ordering, though the gap between graduated and binary (73.1\% vs 72.4\%) is smaller, as perplexity differences attenuate when propagated to task performance.

We scope this as a proof-of-concept with explicit limitations: (1) tested at a single model scale (1B), (2) graduated vs binary confidence intervals partially overlap, and (3) curation effectiveness depends on depth prediction quality.

\subsection{Cross-Scale Probe: 7B Models}
\label{sec:7b}

We probe whether depth signals transfer across model scales by testing DCOL trained on 1--1.5B models against 7B-model recursive chains (Table~\ref{tab:7b}).

\begin{table}[t]
\centering
\small
\caption{Cross-scale transfer to 7B model.}
\label{tab:7b}
\begin{tabular}{@{}lcc@{}}
\toprule
\textbf{Setting} & \textbf{Acc.\ (4-cl)} & \textbf{Acc.\ (6-cl)} \\
\midrule
Zero-shot transfer (1--1.5B $\to$ 7B) & .614 & .372 \\
Fine-tuned on 7B data & .791 & .518 \\
1--1.5B in-distribution (ref.) & .823 & .571 \\
\bottomrule
\end{tabular}
\end{table}

Zero-shot transfer degrades substantially (4-class: 61.4\% vs 82.3\% in-distribution), reflecting that 7B models produce higher-quality generation that shifts the depth boundaries---depth-1 text from a 7B model resembles depth-0 quality at the 1B scale.
However, fine-tuning on 7B data recovers to 79.1\% (4-class), indicating depth signals exist at the larger scale but require recalibration.
EDD at 7B is 3 (nucleus) and 2 (temperature), suggesting that stronger models reach signal saturation at shallower depths.
This motivates future work on scale-invariant depth estimation.

\subsection{Wild-Data Calibration}
\label{sec:wild}

We apply the trained DCOL model to a sample of RedPajama/Dolma web text as an ecological validity probe (Table~\ref{tab:wild}).

\begin{table}[t]
\centering
\small
\caption{Wild-data depth score correlations with quality proxies.}
\label{tab:wild_corr}
\begin{tabular}{@{}lcc@{}}
\toprule
\textbf{Quality proxy} & \textbf{Spearman $\rho$} & \textbf{$p$-value} \\
\midrule
Perplexity (GPT-2) & .341 & .002 \\
AlignScore & .283 & .014 \\
\bottomrule
\end{tabular}
\end{table}

\begin{table}[t]
\centering
\small
\caption{Predicted depth distribution on RedPajama sample.}
\label{tab:wild}
\begin{tabular}{@{}lcccc@{}}
\toprule
\textbf{Depth} & 0 & 1 & 2 & 3+ \\
\midrule
\% of sample & 72 & 16 & 6 & 6 \\
\bottomrule
\end{tabular}
\end{table}

DCOL depth scores show moderate positive correlation with perplexity ($\rho = 0.341$, $p = 0.002$) and weaker correlation with AlignScore ($\rho = 0.283$, $p = 0.014$).
The majority of wild data is scored as depth 0 (72\%), with 28\% receiving depth $\geq$1.
The 6\% scored as depth 3+ may include overestimates, as low-quality human text in the wild can exhibit statistical profiles similar to deeply recursive content.
We present this as a calibration analysis rather than validation, since no ground-truth depth labels exist for wild data.
The controlled benchmark remains our primary evaluation setting.

% ============================================================
\section{Discussion}
\label{sec:discussion}

\paragraph{EDD as a scientific contribution.}
The finding that EDD $\approx$ 3 across model families and decoding strategies provides a concrete answer to a previously open question: \emph{how deep can we look into recursive generation?}
For the model collapse community, this establishes a quantitative boundary on the depth at which statistical signals saturate.
The variation of EDD across conditions (2--4) reveals that model capability and decoding strategy interact to determine this boundary, with weaker models and truncation-based sampling preserving signals longer.

\paragraph{Feature dynamics beyond EDD.}
The Atlas reveals that different features dominate at different depth ranges: surprisal variance is paramount within EDD, but Zipf's $\alpha$ becomes more discriminative beyond it.
This ``feature handoff'' suggests that recursive degradation is not a uniform process but proceeds in phases---first compressing the surprisal distribution, then distorting lexical frequency structure.

\paragraph{Limitations.}
\begin{itemize}[leftmargin=*,topsep=2pt,itemsep=1pt]
\item \textbf{Cross-model gap.} The OLMo rotation fails the 0.45 6-class gate (42.1\%), indicating that cross-model transfer is unreliable when training distributions diverge substantially. Domain adaptation may be needed for deployment across diverse model families.
\item \textbf{Effective granularity.} While we report 6-class results, actual discriminative power is limited to 4 classes (depths 0--3). Depths 4+ collapse into an indistinguishable category.
\item \textbf{Scale transfer.} Zero-shot transfer to 7B models degrades by $\sim$21pp (4-class), requiring scale-specific calibration.
\item \textbf{Controlled vs wild gap.} Our benchmark uses controlled recursive chains; wild synthetic text is produced under diverse, unknown conditions. The moderate wild-data correlations ($\rho \leq 0.34$) indicate this gap is nontrivial.
\item \textbf{DCOL advantage magnitude.} DCOL outperforms the best baseline (CORAL) by +3.3pp. The contribution is real but moderate; much of the ordinal estimation capability comes from the ordinal framework itself rather than DCOL-specific innovations.
\end{itemize}

% ============================================================
\section{Conclusion}
\label{sec:conclusion}

We introduced ordinal generation depth estimation as a new task, moving beyond binary synthetic detection to characterize \emph{how deeply recursive} text content is.
Our Depth Degradation Atlas establishes EDD $\approx$ 3 as the boundary of feasible ordinal discrimination, documents surprisal variance as the most stable depth signal, and maps the interaction between model capability and decoding strategy in determining degradation patterns.
DCOL achieves 82.3\% 4-class accuracy with moderate but consistent gains over baselines.
A proof-of-concept shows that depth-graduated curation outperforms binary filtering, though with modest effect size.
We hope the Atlas---particularly the EDD metric and feature monotonicity analysis---proves directly useful to the model collapse community as a reference for understanding and mitigating recursive contamination.

\begin{thebibliography}{9}
\bibitem{shumailov2024model}
Shumailov, I., et al.
\newblock AI models collapse when trained on recursively generated data.
\newblock {\em Nature}, 631:755--760, 2024.

\bibitem{drayson2025emnlp}
Drayson, H., et al.
\newblock Characterizing non-monotonic degradation in recursive model generation.
\newblock In {\em Proc.\ EMNLP}, 2025.

\bibitem{cao2020coral}
Cao, W., Mirjalili, V., and Raschka, S.
\newblock Rank consistent ordinal regression for neural networks with application to age estimation.
\newblock {\em Pattern Recognition Letters}, 140:325--331, 2020.
\end{thebibliography}

\end{document}
